% ============================================================================
\chapter{Editor and Scene Viewport}
% ============================================================================

\section{Overview}

The scene viewport (\texttt{SceneViewport}) is a 100\% software-rendered
panel driven by ImGui's \texttt{ImDrawList}. There is no GPU projection
matrix; all world-to-screen transformations are implemented analytically
in the \texttt{projectToScreen} function. Three view modes are supported:
\textbf{Mode2D} (orthographic top-down), \textbf{Mode3D} (perspective
orbit), and \textbf{Isometric} (dimetric projection).

\section{Camera Model}

The camera state is stored in \texttt{EditorContext}:

\begin{table}[H]
\centering
\caption{Camera state variables}
\begin{tabular}{lll}
\toprule
\textbf{Variable} & \textbf{Type} & \textbf{Meaning} \\
\midrule
\texttt{camYaw}      & \texttt{f32} & Azimuth angle (radians, Y-axis) \\
\texttt{camPitch}    & \texttt{f32} & Elevation angle (radians, X-axis), clamped to $[-\pi/2, \pi/2]$ \\
\texttt{camFocus}    & \texttt{Vec3} & World-space point the camera orbits \\
\texttt{camDistance} & \texttt{f32} & Distance from focus to eye \\
\bottomrule
\end{tabular}
\end{table}

\section{projectToScreen --- Mode2D}

In 2D mode, the projection is a simple affine pan-and-zoom:

\begin{align}
  s_x &= \text{origin}_x + \frac{W}{2} + p_x \cdot \text{zoom} + \text{panX} \\
  s_y &= \text{origin}_y + \frac{H}{2} - p_y \cdot \text{zoom} + \text{panY}
\end{align}

where $W, H$ are the viewport dimensions and $(p_x, p_y)$ is the world
position.

\section{projectToScreen --- Mode3D}

The 3D projection is a software perspective transform using the
camera's spherical parameterisation.

\subsubsection{Step 1: Translate to Camera-Relative Space}

\begin{equation}
\mathbf{r} = \mathbf{p} - \mathbf{f}
\end{equation}

where $\mathbf{f} = \texttt{camFocus}$.

\subsubsection{Step 2: Yaw Rotation (Y-axis, angle $\psi$)}

\begin{align}
v_x &=  \cos\psi \cdot r_x + \sin\psi \cdot r_z \\
v_y &= r_y \\
v_z &= -\sin\psi \cdot r_x + \cos\psi \cdot r_z
\end{align}

The yaw rotation brings the camera's viewing direction along the $+Z$ axis
in view space.

\subsubsection{Step 3: Pitch Rotation (X-axis, angle $\phi$)}

\begin{align}
v_{y2} &=  \cos\phi \cdot v_y + \sin\phi \cdot v_z \\
v_{z2} &= -\sin\phi \cdot v_y + \cos\phi \cdot v_z
\end{align}

\subsubsection{Step 4: Perspective Divide}

The viewing volume is parameterised by \texttt{camDistance} $D$:

\begin{equation}
\text{fovScale} = \frac{s \cdot D}{\max(D + v_{z2},\; \varepsilon)}
\label{eq:fovscale}
\end{equation}

where $s = \min(W, H) / 2$ is the half-screen size used as a scale
factor, and $\varepsilon = 0.01$ guards against division by zero.
The screen coordinates are:

\begin{align}
  s_x &= c_x + v_x \cdot \text{fovScale} + \text{panX} \\
  s_y &= c_y - v_{y2} \cdot \text{fovScale} + \text{panY}
\end{align}

\subsubsection{Near-Plane Clipping}

A point is \emph{behind the camera} when $D + v_{z2} \leq \varepsilon$.
For line segments (grid, gizmo axes), this causes \texttt{fovScale} to
approach infinity, projecting to extreme screen positions.
The fix is near-plane clipping: for a segment $(\mathbf{a}, \mathbf{b})$,
compute the camera-depth of each endpoint $d_A, d_B$. If exactly one
endpoint is behind the near plane (depth $< \varepsilon$), linearly
interpolate the segment to the boundary:

\begin{equation}
t = \frac{\varepsilon - d_A}{d_B - d_A}, \quad
\mathbf{c} = \mathbf{a} + t(\mathbf{b} - \mathbf{a})
\end{equation}

The clipped segment $(\mathbf{c}, \mathbf{b})$ (or $(\mathbf{a}, \mathbf{c})$)
is then safe to project without numerical explosion.

\begin{specbox}{Invariant 12.1 --- Near-plane clipping}
  Near-plane clipping must be applied to every projected line segment
  in Mode3D. Skipping it for ``short'' segments is not safe; the
  perspective divide can still explode for world-space points near
  the camera position.
\end{specbox}

\section{projectToScreen --- Isometric}

The isometric projection uses a fixed dimetric angle offset of
$\alpha_0 = 30^\circ = \pi/6$ added to the azimuth:

\begin{align}
  \cos A &= \cos(\psi + \alpha_0) \quad\text{(effective azimuth)} \\
  \sin A &= \sin(\psi + \alpha_0) \\
  s_x &= c_x + (r_x \cdot \cos A + r_z \cdot \sin A) \cdot \text{zoom} + \text{panX} \\
  s_y &= c_y - (r_y - r_x \cdot \sin A \cdot 0.5 + r_z \cdot \cos A \cdot 0.5) \cdot \text{zoom} + \text{panY}
\end{align}

This produces the classic isometric look without using a GPU depth buffer.

\section{Infinite Grid Rendering}

The 3D grid is drawn on the XZ plane ($y = 0$). Grid lines are world
segments of the form:

\begin{align}
  &\{(i, 0, z) : z \in [-H, H]\} \quad \text{(X-aligned lines)} \\
  &\{(x, 0, i) : x \in [-H, H]\} \quad \text{(Z-aligned lines)}
\end{align}

where $i \in \{-H_\text{lines}, \ldots, H_\text{lines}\}$ and
$H_\text{lines}$ is chosen dynamically based on \texttt{camDistance}.

Each segment is near-plane clipped before projection. The adaptive
spacing doubles when the grid density exceeds 60 lines on screen:

\begin{equation}
\text{spacing} \leftarrow 2 \cdot \text{spacing}
\quad \text{while } \frac{2 \cdot H_\text{lines}}{\text{spacing}} > 60
\end{equation}

\section{Transform Gizmo}

The \texttt{TransformGizmo} renders translate, rotate, and scale handles
directly on the \texttt{ImDrawList}. For each axis $A \in \{X, Y, Z\}$,
the tip position is computed by projecting the world point
$\mathbf{p} + \Delta A \cdot \hat{\mathbf{e}}_A$ through
\texttt{projectToScreen}.

\subsubsection{Axis Normalisation}

To keep handles at a constant screen-space length $L$ (pixel units),
the raw projected tip is normalised:

\begin{equation}
\text{end}_A = \mathbf{o} + L \cdot \frac{\text{rawEnd}_A - \mathbf{o}}{|\text{rawEnd}_A - \mathbf{o}|}
\end{equation}

If $|\text{rawEnd}_A - \mathbf{o}| < 3\,\text{px}$ (the axis projects
nearly onto the screen-space origin), a per-axis fallback direction is used.

\subsubsection{Z-Axis Overlap Guard}

When the Z-axis fallback direction projects to within
$0.3 \cdot L$ pixels of the Y-axis endpoint, the Z-axis is forced to
its default fallback ($-0.6L, +0.6L$) to avoid ambiguity:

\begin{equation}
d_{ZY} = |\text{end}_Z - \text{end}_Y| < 0.3L
\implies \text{end}_Z \leftarrow \text{fallback}_Z
\end{equation}

\section{Navigation Widget (Orientation Gizmo)}

The mini orientation gizmo in the viewport corner shows the camera's
current viewing direction. Each world axis is projected onto screen space
using the camera rotation only (no perspective divide):

\begin{align}
  s_x    &=  \cos\psi \cdot A_x + \sin\psi \cdot A_z \\
  s_{y1} &= A_y \\
  s_z    &= -\sin\psi \cdot A_x + \cos\psi \cdot A_z \\
  s_{y2} &= \cos\phi \cdot s_{y1} + \sin\phi \cdot s_z \\
  \text{dir} &= \left(\frac{s_x}{L},\; \frac{-s_{y2}}{L}\right) \cdot L_\text{widget}
\end{align}

where $(A_x, A_y, A_z)$ is the world-space direction of the axis (e.g.\
$(1, 0, 0)$ for X), $\psi$ and $\phi$ are \texttt{camYaw} and
\texttt{camPitch}, and $L_\text{widget} = 22\,\text{px}$ is the display
length.

The three axes are drawn in fixed colours: X red (255, 70, 70),
Y green (70, 255, 90), Z blue (90, 140, 255).

\begin{specbox}{Invariant 12.2 --- Navigation widget correctness}
  Navigation widget axis directions must be computed from the live
  \texttt{camYaw} and \texttt{camPitch} values every frame.
  Hardcoding or caching the directions between frames is incorrect.
\end{specbox}

\section{Keyboard Navigation}

Arrow key navigation moves \texttt{camFocus} in the camera's horizontal
plane (pitch is ignored for navigation to avoid tilting the focus point
out of the ground plane):

\begin{align}
  \Delta\mathbf{f}_\text{forward}  &= (-\sin\psi,\; 0,\;  \cos\psi) \cdot v \\
  \Delta\mathbf{f}_\text{right}    &= ( \cos\psi,\; 0,\;  \sin\psi) \cdot v
\end{align}

The \texttt{ImGuiWindowFlags\_NoNavInputs} flag on the viewport window
prevents ImGui from consuming arrow key events to navigate UI buttons.
